\section{来源审计与动机：为什么需要通用感知架构}

本讲讨论的不是“再做一个视觉 Transformer”，而是一个更基础的接口问题：如果输入可能来自图像、声音、触觉、点云、事件相机、文本、蛋白质或科学仪器，我们是否必须为每种数据重新手工设计卷积、tokenizer、局部邻域和专用 decoder？Perceiver 的目标是尽量减少这种逐模态工程，把数据先视为一个带特征与位置的数组，再用统一 attention 接口完成编码、处理和解码。

本次重写以当前 Stanford Online 官方上传 `wTZ3o36lXoQ` 为时间基准。旧讲义引用的 `GV8-6ZgJVRk` 已下线；当前视频提供 1,399 条官方手动英文字幕。录像没有公开独立 slide PDF，因此从 1080p 视频中恢复 39 张完整教学页；重复动画、Q\&A 反复跳页和光流视频中间帧被有意合并，但相应口头解释进入正文与 teacher-voice ledger。

\lecturefigure{slide-01-why-understand-all-modalities.jpg}{为什么要理解所有数据模态：现实世界的数据远不止图像与文本。}{Stanford Online 当前官方视频 00:00:58--00:01:53。}

\begin{knowledgebox}{什么是 inductive bias}
Inductive bias（归纳偏置）是架构预先写入的结构假设。例如卷积假设局部邻域重要且平移共享；tokenizer 假设某种离散子词边界；图神经网络假设输入由节点和边组成。偏置能提高数据效率，却也限定适用数据形态。
\end{knowledgebox}

\teachervoice{Drew Jaegle 的第一个理由非常务实：为每一种传感器和科学数据手工发明偏置，作为研究共同体根本无法持续扩展。通用架构的价值首先是降低“每个新模态一套研究工程”的成本。}

第二个理由是系统维护。传统多模态方案常为视频、音频和文本分别构建 backbone、预处理和 task head，再设计融合模块。它们可以工作，却往往只服务少数固定输入组合，且对采样率、网格形状和模态缺失敏感。Perceiver 希望把架构构建抽象掉，让研究者把精力转向任务、数据和学习目标。

\lecturefigure{slide-02-why-unified-systems.jpg}{从多个专用子系统转向一个可接收异构数组的统一模型。}{Stanford Online 当前官方视频 00:02:58--00:05:09。}

\begin{importantbox}{“通用”不等于“无假设”}
Perceiver 仍假设输入可以表示为有限数组，attention 可以学习元素关系，位置/模态信息可作为特征提供，并且有足够数据恢复任务结构。它减少的是\textbf{模态专用架构假设}，不是消除全部先验。
\end{importantbox}

\begin{warningbox}{统一接口与统一训练是两件事}
讲座展示同一架构分别用于多种任务，但这些实验大多独立训练。能够复用 encode/process/decode 接口，不等于已经训练出一个同时掌握所有模态的单一模型。
\end{warningbox}

\subsection{本章小结}

Perceiver 的问题起点是模态扩张和系统复杂度。它追求可复用的数组接口与较少的手工结构，而不是宣称领域知识、专用模型或联合多模态训练已不再需要。

